\section{Shock fitting and the postshock flow}
\label{sec:shock}
\subsection{A fitted shock is an unknown boundary}
Shock fitting represents a discontinuity explicitly and solves for its
position/orientation together with the flow on each side. Shock capturing
instead represents the jump through a discretized conservation-law field.
The original MOC treatment matches downstream characteristic information
to the shock relations \cite[\S11.5.3, PDF p.~411;
\S11.15, PDF p.~438]{anderson}.

The current code's cell-fold or coalescence detection is a useful refusal
mechanism for its smooth march. A fold alone does not select an admissible
shock, establish its location, or prove that a physical shock rather than
under-resolution caused the failure. No fitted-shock solve was located in
the current \code{raosim} paths.

\subsection{Conservation is necessary but does not locate the front}
For a locally stationary adiabatic inviscid shock without external work, with normal directed from
upstream 1 to downstream 2 and nonzero normal mass flux \(j\),
\begin{align}
j&=\rho_1V_{n1}=\rho_2V_{n2},\\
p_1+\rho_1V_{n1}^2&=p_2+\rho_2V_{n2}^2,\\
V_{t2}&=V_{t1},\\
h_1+\tfrac12|\mathbf V_1|^2&=
h_2+\tfrac12|\mathbf V_2|^2 .
\label{eq:RH}
\end{align}
These are local conservation relations, with tangential continuity for
the stated inviscid stationary setting. They are not restricted to a
calorically perfect gas
\cite[\S3.6, PDF p.~104, Eqs.~3.38--3.40;
\S4.3, PDF pp.~150--151, Eqs.~4.3--4.6]{anderson}.

They must be closed using the same EOS, enthalpy/composition convention
and acoustic mode as the rest of the flow. Upstream \(p,T,\mathbf V,Y\),
the front normal, and an entropy-admissible compressive root are needed.
The trivial unshocked root must be excluded. For a general reacting EOS,
an upstream Mach number alone does not specify the dimensional state
\cite[\S17.2, PDF pp.~665--668]{anderson}.

Front geometry is determined by additional characteristic, wall,
symmetry-axis and downstream-pressure conditions. The familiar
\(\theta\)--\(\beta\)--\(M\) relation is a calorically perfect limiting
test, not a general-gas closure with locally substituted \(\gamma\).
Weak/strong roots, attached-front limits and pressure-dependent branch
selection must remain explicit \cite[PDF pp.~152--154]{anderson}.

\subsection{Curved shocks require rotational transport}
A curved front generally produces different entropy jumps between
streamlines when its shock strength varies.
The postshock state therefore needs more information than the current
smooth isentropic C+/C- invariants. Anderson's rotational MOC carries
streamline entropy and total enthalpy in addition to the Mach-line data,
including an interpolated streamline foot
\cite[\S11.9, PDF pp.~423--424]{anderson}.

\proposal Represent separate flow regions and an explicit shock-front
mesh. Shock nodes should reference their upstream sampling location,
front normal, jump result and downstream characteristic parents.
Transport entropy/enthalpy or equivalent conservative thermodynamic
information into the downstream rotational unit processes. Include
composition and relaxation history for a reacting branch.

Conserve total enthalpy rather than assuming unchanged total temperature
for all gas models. Likewise, the calorically perfect relation between
entropy rise and stagnation-pressure loss cannot be transferred without
qualification to changing-composition states
\cite[PDF p.~109, Eqs.~3.61--3.63]{anderson}.

\subsection{Frozen fronts, reacting relaxation and mixed-type limits}
Frozen jump conditions followed by finite-rate downstream relaxation
describe a different model from an equilibrium jump. The source
distinguishes the very short shock-front scale from chemical relaxation
downstream \cite[\S17.10, PDF pp.~690--691]{anderson}.
The selected model should name these assumptions and their reference
evidence, rather than silently relaxing to a preferred equilibrium state.

Steady supersonic MOC requires a locally supersonic marching region.
A normal shock or sufficiently strong oblique shock can create a
subsonic downstream region. Separation also introduces recirculation and
viscous interaction. These regions require a coupled mixed-type or
viscous method with the appropriate boundary conditions. A hyperbolic
characteristic march alone cannot supply that solution
\cite[PDF p.~411]{anderson}\cite[PDF pp.~4--5]{pns}.

The first fitted-shock milestone should therefore be deliberately
restricted: a specified stationary topology and gas mode, with an
admissible supersonic downstream region where the chosen MOC applies.
Shock--wall/axis intersections, shock interactions and slip lines are
additional topology milestones. General off-design back-pressure
analysis needs an explicit subsonic/separated-flow route.

\section{Viscous interaction, heat transfer and separation}
\label{sec:viscous}
\subsection{Attached boundary layers are a separate flow model}
\current The repository has heat-transfer/cooling models, a
displacement/area-loss screen and empirical separation criteria.
\code{physics.boundary_layer_displacement} labels its nozzle extension
as a heuristic without a located validity range. None evolves a
compressible nozzle boundary layer and iterates its displacement into MOC.

NASA SP-8120 calls for integral momentum and energy boundary-layer
analysis with friction/heat-transfer closures, and for displacement
thickness in performance prediction \cite[PDF p.~27, printed p.~13]{sp8120}.
That passage supplies the modeling requirement and references; it is not
itself a complete axisymmetric closure set ready to transcribe.

\proposal Introduce an attached boundary-layer state with source-selected
definitions for momentum/energy thickness or profiles; edge
\(p_e,\rho_e,U_e,h_e,Y_e\); transition/relaminarization assumptions;
wall temperature/chemistry/roughness; skin friction; heat flux; and
pressure-gradient/axisymmetric terms. Transport and thermal property
models need independent identity and applicability.

A proposed local displacement convention, derived by replacing the
integrated mass deficit with an equivalent edge-state strip and written
in a wall-normal coordinate \(n\), is
\begin{equation}
\delta^*=\int_0^\infty
\left(1-\frac{\rho u}{\rho_eU_e}\right)\,\dd n .
\end{equation}
This is a project definition of the local mass deficit, not a
source-transcribed closure. It is not the complete
axisymmetric integral boundary-layer method or a proposed universal
velocity-profile closure. Those geometry and closure choices remain
implementation dependencies.

\subsection{Preserve the physical wall and distinguish forward from inverse use}
For a prescribed metal wall, displacement reduces the effective core.
In a thin, locally near-axial approximation one may interpret
\begin{equation}
A_{\mathrm{core}}\simeq\pi(r_w-\delta^*)^2.
\end{equation}
This is a proposed geometric interpretation for forward analysis.
The actual implementation must derive the correct local-normal and
axisymmetric projection rather than applying this scalar radial formula
to an arbitrary steep wall.

SP-8120's outward wall correction is an inverse-design compensation:
move the design wall outward to retain a desired inviscid core
\cite[PDF p.~27]{sp8120}. A forward analysis call should not alter the
manufactured wall or export displacement as hardware geometry.
Maintain physical-wall, core-edge and thermal-station coordinates
separately and iterate their consistent coupling.

\subsection{Energy and loss accounting must close with the coolant solution}
Skin friction changes momentum and total-pressure loss. Wall heat loss
changes gas enthalpy, coolant enthalpy and performance. A local heat
correlation used by the cooling module is not a resolved viscous flow field.
A coupled calculation should pass the edge state and wall temperature
between core/BL/thermal models and close a common energy ledger.

The Wang nozzle study explicitly couples species, momentum and total
enthalpy, and compares adiabatic and cooled-wall predictions
\cite[PDF pp.~2, 5, 7]{wang}. Its mechanism and radiation observations
are case-specific; they do not establish a universal mechanism or justify
omitting radiation in every propellant.

When a physical loss is already represented in the flow and exit flux,
do not subtract the same effect again through a drag decrement and a
generic \(\eta_{CF}\). Residual efficiency factors may remain for
explicitly unmodeled losses, with their meaning recorded.

\subsection{Thin attached-layer corrections have an applicability boundary}
Strong acceleration, transition, relaminarization, wall cooling and
low Reynolds number complicate nozzle boundary layers. When displacement
is comparable to the nozzle radius, a separate inviscid core plus small
correction may break down. Streamwise-separated regions also violate the
assumptions of simple downstream viscous marching
\cite[PDF pp.~4--5]{pns}\cite[PDF p.~33]{sp8120}.

No universal displacement ratio or Reynolds cutoff is adopted here.
Source-selected limits and validation-domain decisions must accompany
the model. A residual converging numerically does not establish that
its attached-layer assumptions remain applicable.

\subsection{Separation onset, pattern, reattachment and side loads}
\current \code{separation.py} implements Summerfield, Schilling,
Kalt--Badal and Schmucker onset screens, with explicit FSS-only scope.
The local MDO uses a Schmucker margin. Wall-station inputs are reconstructed
from Q1D area states or the separate frozen expansion; the function does
not consume a resolved MOC pressure/BL field.

The original separation review discusses scatter and the distinction
between shock-interaction pressure rise and recirculation recovery
\cite[\S6.3.1, PDF pp.~51--52]{ostlund}. SP-8120 describes separation
criteria as development guides and calls for full-scale checks of
marginal designs \cite[PDF p.~82]{sp8120}. Existing screening should be
retained with its scope, not relabeled as resolved separated flow.

A broader suite must distinguish:
\begin{itemize}
\item incipient interaction and actual separation location;
\item free shock separation (FSS), restricted shock separation (RSS),
reattachment and recirculation/plateau pressure;
\item steady branch selection versus startup/shutdown history;
\item mean axisymmetric loads versus unsteady asymmetric side loads.
\end{itemize}
For the cited VOLVO S1 case, the source reports different FSS/RSS states
at the same pressure ratio on startup and shutdown, with transitions
near 15 and 12 respectively \cite[\S8.1.1, PDF p.~70]{ostlund}.
These are case observations, not transferable default thresholds.
A memoryless smooth NPR-to-pattern function cannot reproduce both states.

For the pressure contribution of a circumferentially symmetric solution,
\[
\mathbf F_\perp=\int_{\mathrm{wall}}p\,\mathbf n_\perp\,\dd A=\mathbf0
\]
by symmetry; symmetric viscous tractions also cancel laterally.
A steady axisymmetric MOC/BL result therefore cannot itself
provide a directional side-load time history. Side-load prediction needs
an independently calibrated statistical model or asymmetric unsteady
viscous simulation and structural/measurement transfer information
\cite[PDF pp.~70, 78--82]{ostlund}\cite[PDF p.~6]{frey}.

\proposal First connect source-appropriate onset screens to resolved wall
edge states while retaining their calibrated Mach, geometry and thermal
scope. Then develop a separately identified shock/BL interaction and
separated mixed-type model. RSS and side-load products require their own
validation, even when offered behind a common user-facing solver suite.
